\documentclass[11pt]{amsart}
\usepackage{amsmath,amssymb,amsthm}
\usepackage[margin=1.1in]{geometry}
\usepackage{newtxtext,newtxmath}
\usepackage{microtype}
\usepackage{booktabs}
\usepackage[colorlinks=true, linkcolor=blue, citecolor=blue, urlcolor=blue]{hyperref}
\usepackage{listings}
\usepackage{upquote}
\lstdefinelanguage{Lean}{
  morekeywords={def,theorem,abbrev,open,namespace,end},
  sensitive=true,
  morecomment=[l]{--},
}
\lstset{language=Lean, basicstyle=\small\fontfamily{lmtt}\selectfont\upshape, keywordstyle=\bfseries,
  columns=fullflexible, keepspaces=true, breaklines=false,
  xleftmargin=1em, aboveskip=6pt, belowskip=6pt,
  literate={ℕ}{{$\mathbb{N}$}}1 {ℝ}{{$\mathbb{R}$}}1 {ℤ}{{$\mathbb{Z}$}}1
    {→}{{$\rightarrow$}}1 {×}{{$\times$}}1 {∃}{{$\exists$}}1 {∀}{{$\forall$}}1
    {κ}{{$\kappa$}}1 {≤}{{$\le$}}1 {≥}{{$\ge$}}1 {∧}{{$\wedge$}}1
    {¬}{{$\neg$}}1 {∈}{{$\in$}}1 {₀}{{$_0$}}1 {₁}{{$_1$}}1 {₂}{{$_2$}}1 {⊕}{{$\oplus$}}1
    {ᶠ}{{$^{f}$}}1 {ℓ}{{$\ell$}}1
}

\newtheorem{theorem}{Theorem}
\newtheorem{lemma}[theorem]{Lemma}
\theoremstyle{remark}
\newtheorem{remark}[theorem]{Remark}

\begin{document}

\title[A fast Mathlib multiplier]{Integer multiplication below $n\log n$\\ on Mathlib's Turing machines}

\author{Dan Brown$^*$}
\email{jdb1729@gmail.com}
\thanks{$^*$Author of \emph{The Da Vinci Code}.}

\author{Kenny Daniel}
\email{platypii@gmail.com}

\date{October 14, 2026}

\begin{abstract}
We prove in Lean, against Mathlib's \texttt{Turing.FinTM2} and \texttt{TM2ComputableInTime},
that a fixed finite multi-stack machine multiplies two naturals in worst-case time
$O(n(\log n)^{1-\kappa})$ for some $\kappa>0$. The machine is one host program whose stages are
the numeric algorithm of the paper: a CRT reindexing, two linewise resamplings, chirp
products, a complex convolution through computed synthetic transforms whose rounds are the
simultaneous butterfly layer of \S5, pointwise products, and a final rounding. Every stage is
a verified machine with an explicit step bound, and the sum of the bounds is
$O(n(\log n)^{1-\kappa})$ with $\kappa=\varepsilon c(1-\tau)$ explicit. The vendored library
of Kenny Daniel supplies the network, the interchange, the Schönhage--Strassen products and
the resampling mathematics. The development has 248 modules and 134\,932 lines; every main
theorem depends only on \texttt{propext}, \texttt{Classical.choice} and \texttt{Quot.sound}.
\end{abstract}

\maketitle

\section{The theorem}

The statement (\texttt{Statement.lean}) is Mathlib vocabulary. The input $(a,b)$ is
\texttt{encodingProd} of two copies of \texttt{encodingNatBool}: the LSB-first binary word of
$a$ over \texttt{Sum.inl} symbols followed by that of $b$ over \texttt{Sum.inr}; its length is
$n=\texttt{Nat.size}\,a+\texttt{Nat.size}\,b$. The output is the canonical LSB-first word of
$ab$. The time function of \texttt{TM2ComputableInTime} bounds the step count on every input
of length $n$, so the bound is worst case.

\begin{lstlisting}
abbrev pairEnc : ℕ × ℕ → List (Bool ⊕ Bool) :=
  (encodingProd encodingNatBool encodingNatBool).encode

def MulInTime (κ : ℝ) : Prop :=
  ∃ M : TM2ComputableInTime pairEnc encodingNatBool.encode
          (fun p : ℕ × ℕ => p.1 * p.2),
    (fun n : ℕ => (M.time n : ℝ)) =O[atTop]
      fun n : ℕ => (n : ℝ) * Real.log n ^ (1 - κ)

def Statement : Prop := ∃ κ : ℝ, 0 < κ ∧ MulInTime κ
\end{lstlisting}

\begin{theorem}[\texttt{Theorem.lean}]
\texttt{theorem statement : Statement}.
\end{theorem}

The proof is one line: \texttt{Fast/Final.fast\_statement\_of} applied to
\texttt{Cost/MatchSynLoop.mr\_loop\_cost\_le}. The first is the host program with every bundle
instantiated; it takes as its only hypothesis a bound on the cost of the round loops of the
synthetic transforms, and the second proves that bound. \texttt{FastPathCheck.lean} asserts with
\texttt{\#guard\_msgs} that \texttt{statement} and the main intermediate theorems depend on
\texttt{propext}, \texttt{Classical.choice} and \texttt{Quot.sound} only. The exponent is
$\kappa=\varepsilon c(1-\tau)=296/(10^{11}\cdot 2^{131})$ (\texttt{Cost/Total.tot\_kappa\_eq}).

\section{The reduction to one program}

A \emph{host program} is a procedure table in the structured stack language over the Mathlib
input and output stacks and a bank of register stacks. \texttt{HostMulWithin body pm v$_0$ time x
y} says that from the store holding $x\#y$ on the input register the procedure \texttt{pm}
halts in the exact final configuration required by \texttt{TM2ComputableInTime}, with the
product on the output register, in at most $\texttt{time}(n)-2$ steps.

\begin{theorem}[\texttt{Fast/Glue.statement\_of\_dispatch}]
If a host program satisfies \texttt{HostMulWithin} on every input of encoded length at least
$n_0$ and its time is $O(n(\log n)^{1-\kappa})$ with $\kappa>0$, then \texttt{Statement} holds.
\end{theorem}

A dispatcher reads up to $n_0$ symbols into its finite control, answers a short input from a
finite table, and otherwise pushes the symbols back and runs the program.
\texttt{Cost/Total.tot\_htime} turns any eventual bound $\texttt{time}(n)\le\texttt{totBound}\,C\,n$
on an $\mathbb N$-valued time into the hypothesis \texttt{htime} at $\kappa=\texttt{totKappa}$.

\section{The numeric algorithm}

\texttt{Mult/Pipeline.pipelineCoeffs n x y} is the computation of the paper with every
ingredient fixed: the parameters $b,p,d,K,T,r,\ell,t_i,s_i,\alpha,\gamma$ of
\texttt{Mult/Sizes}; the CRT reindexing $\Phi$ of \texttt{Mult/CRT}; the line maps
$\widetilde A$ and $\widetilde B_0$ with the vendored weight tables; the chirp table; the
computed complex convolution through twisted synthetic transforms; pointwise products; the
final scaling and rounding.

\begin{lstlisting}
theorem pipeline_correct :
    ∀ᶠ n in Filter.atTop, ∀ x y : ℕ, Nat.size x ≤ n → Nat.size y ≤ n →
      evalBase (2 ^ b n) (pipelineCoeffs n x y) = x * y
\end{lstlisting}

The moduli are pairwise coprime (\texttt{Mult/Moduli.src\_good}): group A moduli $Nx+1$ and
group B moduli $Ny+2$ with $N$ the product of the odd numbers below $4d$, so no prime number
theorem is used. The source transform is within $9\cdot 2^\gamma T\log_2T$ grid units of the
exact one (\texttt{Mult/SourceErr}), the synthetic convolution error is
$\sqrt2\,M(3\ell+1)$ (\texttt{Mult/SynthErr}), and the forward and opposite synthetic
transforms are exactly $\ell$ rounds of normalized butterflies on one selected bit per axis
followed by per-record twiddles $y^{\mp E}$ (\texttt{Mult/Rounds}). $\widetilde B_0$ runs
$\alpha^2$ Neumann steps on a window of $\lfloor\sqrt p\rfloor/(2\alpha)+1$ terms
(\texttt{Mult/FastB0}). Fixed-point products inside the per-entry stages go through the
vendored Schönhage--Strassen multiplier.

\section{The host program}

\texttt{Fast/MainProg.mnFast} is the single host program. Its register file
(\texttt{Fast/MainRegs}) is $\texttt{MnAll}=\texttt{MnR}\oplus\texttt{Fin 64}\oplus
\texttt{LReg SXReg}$: the stage registers, the 64 registers of the Schönhage--Strassen
multiplier, and a reserved copy of the register file of the layer machine for the synthetic
transforms and the parameter distribution. The stages, with the cost term of each in
\texttt{Cost/Total.totBound}:

\begin{center}\small
\begin{tabular}{@{}llll@{}}
\toprule
& stage & semantics & cost term\\
\midrule
F1 & read $x$, $y$ & \texttt{readInput\_exec} & $n$\\
F2 & parameters, tables & \texttt{pmParams\_exec}, table machines & $n+B^C$, $Tp$\\
F3 & scaled digit records & \texttt{digitRecords\_exec} & $n$\\
F4 & $\Phi$: stream to padded box & \texttt{crtForward\_runs} & $Tp\,d^2(1+\ell^\tau)$\\
F5, F6 & $F'$ on $x$ and on $y$ & \texttt{srcP\_runs} & resampling, convolution, $TM(p)$\\
F7 & pointwise products & \texttt{cst\_point\_runs} & $TM(p)$\\
F8 & $G'=\overline{F'(\bar\cdot)}$ & \texttt{srcOpp\_runs} & as F5\\
F9 & $\Phi^{-1}$: box to stream & \texttt{crtInverse\_runs} & $Tp\,d^2(1+\ell^\tau)$\\
F10 & rounding & \texttt{roundStage\_exec} & $S\,M(p)$\\
F11 & carries, output & \texttt{carryOut\_exec} & $n$\\
\bottomrule
\end{tabular}
\end{center}

Here $B=b(n)=\lceil\log_2n\rceil$, $M(p)=p^{1+\delta}$ with $\delta=1/8$ is the cost of one
$p$-bit product, the resampling term is $d\,T\,p^{3/2+\delta}\alpha$, and the convolution term
is $Tp\,(dK+pK^{\tau-1}+\ell d^{\lambda'}+\log(rp)\log\log(rp))$. The program ends by emptying
every register and resetting its local state.

Two parts of \texttt{mnFast} are abstract parameters with contracts: the parameter
distribution (\texttt{MnDistSpec}) and the synthetic transforms (\texttt{SynFwdSpec},
\texttt{SynOppSpec}, \texttt{Fast/ConvSpec}). \texttt{Fast/MainCorrect} proves correctness beyond
the threshold $n_0$ and \texttt{Fast/MainTime} the bound $\texttt{mnTime}\le\texttt{totBound}\,C$.

\begin{theorem}[\texttt{Fast/MainStatement.statement\_of\_specs}]
If for some coefficient width $w$ there are a layout bundle $L:\texttt{MnLay}\ w$ and a
distribution bundle \texttt{MnDist mnKa mnLb w (mn$\sigma$L L)}, then \texttt{Statement} holds.
\end{theorem}

\texttt{Fast/BundleMnLay.blLay} builds \texttt{MnLay} from the verified layout programs
\texttt{slyFwd}, \texttt{slyOpp} (\texttt{Fast/SynLayFwd}, \texttt{Fast/SynLayOpp}), the
parameter wires $v_0$ and the transform cost bound. \texttt{Fast/BundleDist.bdDist} builds
\texttt{MnDist} from the verified distribution machine \texttt{dsMain}, the rounding parameters
(\texttt{Fast/BundleRnd}) and a wires program that writes the transform parameters on the
reserved layer registers; the wires program is \texttt{Fast/DistSynMain.dsyProg}
(\texttt{Fast/BundleWires}). \texttt{Fast/MainFinal.statement\_of\_wires} needs only the layout
bundle and the wires program, and \texttt{Fast/Final.fast\_statement\_of} supplies both from
the loop-cost bound.

\section{The stage machines}

Every machine comes with an \texttt{\_exec} or \texttt{\_runs} theorem giving its exact final
register contents and an explicit step bound.

\begin{itemize}
\item \emph{Input and output} (\texttt{Fast/Input}, \texttt{Fast/Output}): the operands LSB
first, the unary input length, the digit records $\mathrm{digit}\cdot 2^{p-b-2}$ zero padded
to $S$; carry propagation in radix $2^b$ and the Mathlib-encoded output.
\item \emph{Parameters} (\texttt{Fast/Param*}): a binary table machine with 40 variables of
$4\,\texttt{size}\,n+16$ bits computes every parameter of \texttt{Mult/Sizes},
\texttt{Mult/Moduli}, \texttt{Mult/Pipeline} and \texttt{Mult/Chirp} exactly from the unary
input length, emits the axis lengths as word streams, and pushes unary copies
(\texttt{pmParams\_exec}, cost $16n+O((\texttt{size}\,n)^4)$ plus the unary copies).
\item \emph{Tables} (\texttt{Fast/Table*}): the same register machine
(\texttt{Fast/TableVM}, at most $32(W+4)^2+8w$ host steps per atom) computes Machin's $\pi$
(\texttt{tblPi\_eq\_piQ}), the Taylor exponentials and the weight tables of $\widetilde A$ and
$\widetilde B_0$, the chirp axis tables and the twist tables, as streams of Gaussian numerators
in the order the pipeline consumes them.
\item \emph{Distribution} (\texttt{Fast/Dist*}): one fixed program after the parameter machine
writes the record and stream formats read by every stage loader (the per-axis records of the
resamplings, the chirp blocks, the layout parameters, the ring header, the layer constants)
and the register file between runs of the source transform
(\texttt{dsAll\_exec}, \texttt{dsAll\_base}; cost $2\cdot10^{10}(n+(b+1)^5)$).
\item \emph{CRT reindexing} (\texttt{Fast/CRT*}): per target axis a walk moving the axis below
the preceding ones and a rotation pass; forward with a padding scan, inverse with a deletion
scan (\texttt{crtForward\_runs}, \texttt{crtInverse\_runs}, round trip
\texttt{crtRoundTrip\_runs}).
\item \emph{Resampling} (\texttt{Fast/Prep*}, \texttt{Fast/Resamp*}): the line machines of
$\widetilde A$ and of the fast $\widetilde B_0$ applied along every axis of the padded box in
the order of \texttt{tensorZR}, the last axis directly and every other axis through the
vendored interchange with the last one (\texttt{rsaStage\_exec}, \texttt{rsbStage\_exec};
\texttt{rsa\_all\_lineA}, \texttt{rsb\_all\_tensor}).
\item \emph{Chirp and pointwise} (\texttt{Fast/Chirp*}, \texttt{Fast/Point*},
\texttt{Fast/CMul*}): the chirp kernel box from the axis tables by repetition passes and table
products (\texttt{cst\_chirpBox\_kernel}); entrywise chirp products, conjugation, the exact
scale by $2^\gamma$ and the pointwise product (\texttt{cst\_point\_runs}), each through the
complex multiplier \texttt{cmulStream} over Schönhage--Strassen.
\item \emph{Convolution} (\texttt{Fast/Conv*}, \texttt{Fast/Ring*}, \texttt{Fast/Load*}): the
twist, the synthetic convolution and the untwist on box streams
(\texttt{cvBox\_convBox}, the instance \texttt{Pipeline.convBox n}); the pointwise negacyclic
ring products of two record streams by one placed run of the vendored signed ring product
(\texttt{rpProg\_exec}); the loaders of the chirp kernel, the layout stage and the ring header
(\texttt{ldChirpStage\_exec}, \texttt{ldLayGenStage\_exec}, \texttt{ldRpStage\_exec}).
\item \emph{Rounding} (\texttt{Fast/Round}): coefficient $i$ is
$\mathrm{round}(2^{2b+4}S^2\,\mathrm{re}\,N_i/2^p)$, computed by a sign mask, one
Schönhage--Strassen product and a shift (\texttt{roundStage\_exec}, \texttt{rndSquare\_exec}).
\end{itemize}

\section{The synthetic transforms and the layer machine}

\subsection*{The layer machine}
\texttt{Node/*} is the machine of \S5. One procedure \texttt{layer} on the register file
\texttt{LReg} (the array \texttt{io}, the bank of wire registers with temps, a LIFO parking
register for the parent's wires during a child call, gate and shift counters, phase, leaf,
basis-change and route registers, unary parameter registers) computes the recursive
\texttt{layerZ} on Gaussian integers modulo $2^w$: at a node the array is split into roles,
the input phases are applied, the instruction list of the vendored network runs with
\texttt{gateProg} for gates and \texttt{factorProg} for factors (a basis change by slot XORs,
a child call with the other wires parked, the basis change undone), the output phases and the
unit $i^{27k}$ are applied, and the roles merge. \texttt{Node/LayerProc.layer\_exec} is proved
by strong induction on the level. \texttt{Node/TopProc.topProg} is one fixed program computing
\texttt{topZ} (phase $S$, kernels on the top chunks, rows and padding, the base-$m$ pieces
through \texttt{layer}, phase $S$, the unit $(1-i)^D2^{-\lfloor D/2\rfloor}$, the truncating
shift); every size and configuration value is read from registers at run time, the pieces from
a record consumed by a loop (\texttt{topProg\_exec}).

\subsection*{The synthetic transforms}
\texttt{Fast/Syn*} runs the rounds of \texttt{Mult/Rounds} on the layer machine. The records
stay in the physical order of the layout $L=\texttt{layPerm}\;es\;\ell\;K$, written by the
layout programs. For each round $h$ the program \texttt{synRec} computes the configuration and
piece records of \texttt{topProg} from unary parameters (\texttt{Fast/SynRec}). A chunk round
selects the bits $\rho+K(c_0+j)$, $j<D$; the highest
\[
 q_0=\min\bigl(46\lfloor\log_mD\rfloor,\,D\bigr)
\]
of them are top kernels, the rows are the top $a+Kq_0$ address bits, and the remaining
$D'=D-q_0$ bits are selected by one piece per nonzero base-$m$ digit of $D'$ on
$R_2=2^X\,NW^{J'}$ rows with $2^{a'}\le R_2<2^{a'+1}$: the padding at most doubles the volume.
The forward and opposite transforms are one loop over the round counter
(\texttt{Fast/SynFwd.synFwd\_spec\_br}, \texttt{synOpp\_spec\_br}), with the twiddles of
\texttt{Fast/Twiddle*}.

\subsection*{Register embedding}
The transforms are proved on the register file of the layer machine and run inside the
multiplier's file. \texttt{Host/Embed.exec\_emb}: an execution of $p$ in $n$ steps on $T$
tapes gives an execution of \texttt{embProg e $\iota$ p} in exactly $n$ steps on $T'$ tapes along an
injective $e:\texttt{Fin}\,T\to\texttt{Fin}\,T'$, from every big store that reads back to the
small one, ending in a store that reads back to the small final one and is unchanged off the
image of $e$. \texttt{Fast/SynFrame.embSynFwd\_spec} restates both contracts in the big file.

\section{The cost}

\subsection*{The analytic bound}
\texttt{Cost/Total.totBound C n} is $C$ times the sum of the stage terms of the table above,
written in the parameters of \texttt{Mult/Sizes}. The exponents in force are the vendored
$\tau=1-296/10^{11}$ (interchange) and $\sigma=1-418/10^{12}$ (branching),
$\beta=1/2$, $\lambda=1-3/10^{10}$, $\lambda'=1-1/10^{10}$, $\delta=1/8$, and
$\varepsilon=2^{-75}$, $c=2^{-56}$ from the sizes. With $Tp\le48n$ each term is at most a
constant times $nB^x$; the largest exponent is $1-\kappa$, from the convolution's
$Tp\cdot pK^{\tau-1}$, since $B^{\varepsilon c}\le 4K$. The ring-product term
$Tp\,\log(rp)\log\log(rp)$ has exponent $1-\varepsilon+\varepsilon c$.

\begin{lstlisting}
theorem totBound_isBigO (C : ℕ) :
    ∃ κ : ℝ, 0 < κ ∧
      (fun n => totBound C n) =O[atTop]
        fun n : ℕ => (n : ℝ) * Real.log n ^ (1 - κ)
\end{lstlisting}

\texttt{Cost/Growth} holds the bounds on the parameters in powers of $B$ ($d\le B^\varepsilon\le
2d$, $K\le B^{\varepsilon c}$, $24n\le Tp\le 48n$, $d\ell\le 2B$, $\alpha\le 4d^{1/2}B^{1/4}$)
and the closure lemmas for \texttt{IsBigO} against $n(\log n)^e$, including the
Schönhage--Strassen step bound $m(\log m+1)(\log\log m+1)\le 625\,m^{9/8}$.

\subsection*{The layer recurrence}
\texttt{Cost/Branching} computes the factor count $F$ of the vendored network's program as the
rank sum of its edges and proves $F/NW<m^\sigma$ with $NW=58645352620000$.
\texttt{Cost/LayerCost} puts the step bound of \texttt{layer} in per-volume form:
$G(0)\le 200u(1+n/(rw))$ and
$G(j+1)\le (F/NW)\,G(j)+C_1(1+\epsilon)((Km^ju)^\tau+1)$ with $\epsilon\le1$ at admissible
nodes. \texttt{Layer/Recurrence.layer\_cost\_bound} solves the recurrence: under
$\max(\tau,\sigma)<\lambda<\lambda'$, $\tau(1+c/\beta)\le\lambda$ and
$\sigma+\beta(1-\sigma)<\lambda'$, the base-$m$ pieces of $D'\le d$ chunks plus
$q_0\le Q(1+\log_md)$ kernels cost at most $C\,d^{\lambda'}$ normalized, with one constant for
all $d$, $K$, cutoffs and piece decompositions. \texttt{Cost/Total.tot\_lam\_conditions}
checks the inequalities at the exponents in force.

\subsection*{Matching the machines}
\texttt{Cost/Match*} bounds every stage machine by its term, eventually, with constants
depending only on the fixed data (the register count, the interchange bank size, the alphabet,
the vendored constants, the width ratios $w\le c_wp$, $W\le c_Wp^2$). Two charges go beyond
the intended terms. The resamplings interchange each axis with the last one, $2(d-1)$
interchanges of volume $3Tw$ at $\Theta(dTp\ell^\tau)$ in all; this exceeds the resampling
term and is charged to the CRT term $Tp\,d^2(1+\ell^\tau)$. The layout writer's table word has
$rw+N^2+N+4$ bits, about $r\cdot w$, and the writer costs a polynomial in $B$ and $r$, inside
$Tp$ (\texttt{ml\_poly\_r\_le\_T}). The synthetic transforms are matched in three parts: the
layout programs against $Tp(dK+pK^{\tau-1})$ (\texttt{Cost/MatchSynLay}), the layer of a chunk
round against $Tp\,d^{\lambda'}$ through \texttt{layer\_cost\_bound}
(\texttt{Cost/MatchSynPcs}, \texttt{MatchSynTop}), and the records, rest states, parking and
twiddles against $Tp$ per round (\texttt{Cost/MatchSynOver}).

\begin{theorem}[\texttt{Cost/MatchSynLoop.mr\_loop\_cost\_le}]
For every $C_s$ and $k$ there is $C$ such that for all large $n$ and every parking cost
$\mathrm{stk}\le C_s(r+(b+1)^k)$, the forward and opposite round loops of the synthetic
transforms at the pipeline parameters cost at most $C\cdot\texttt{totConv}\,n$ steps.
\end{theorem}

\section{The vendored library}

The directory \texttt{IntegerMultBounds/} is the library of Kenny Daniel (MIT license). The
development uses its network (\texttt{Networks/}: the compact complex network with
$W=58645352620000$ wires and label dimension $m=25^3$, its directions and the branching bound
\texttt{ComplexRank25.branching\_bound}), its interchange of address chunks
(\texttt{Swap/}, \texttt{BinaryInterchangeRun.uniform\_bound}, cost $CV\max(1,L)^\tau$), its
Schönhage--Strassen multiplier and signed ring product (\texttt{Schoenhage/}), its resampling
mathematics (\texttt{Resampling/}: the weight tables, the Neumann series of $\widetilde B_0$,
the operator-norm bound \texttt{opNorm\_offDiagCLM\_le\_rpow}), and \texttt{ExactRecovery}.
Each vendored program runs inside the host as a foreign call with a \texttt{HoareTime}
contract on host registers (\texttt{Host/Core}, \texttt{Host/Placed}).

\section{Size and verification}

\texttt{IntMul.lean} imports 248 modules, 134\,932 lines in all: \texttt{Fast} (144 modules,
78\,548 lines), \texttt{Node} (20, 21\,268), \texttt{Cost} (21, 9\,054), \texttt{Host}
(18, 8\,732), \texttt{Mult} (16, 8\,234), \texttt{Layer} (19, 6\,442), \texttt{Bridge}
(7, 2\,085), \texttt{Machine} (1, 512), and \texttt{Statement.lean}, \texttt{Theorem.lean}.
\texttt{lake build} of the library passes. \texttt{FastPathCheck.lean} asserts with
\texttt{\#guard\_msgs} that \texttt{statement} and the main intermediate theorems depend on
\texttt{propext}, \texttt{Classical.choice} and \texttt{Quot.sound} only, and
\texttt{MathlibTM2Check.lean} does the same for the bridge theorems of \texttt{Bridge/Main}
(\texttt{mul\_quadratic}, \texttt{statement\_of\_endToEnd}, \texttt{tm2\_of\_runsWithin},
\texttt{mulInTime\_of\_computesWithin}); CI runs both.

\end{document}
